\documentclass{article}
\usepackage{amsmath,amssymb,amsthm}
\usepackage{algorithm}
\usepackage[noend]{algpseudocode}
\usepackage{bm}
\usepackage{enumitem}
\usepackage{hyperref}
\newtheorem{theorem}{Theorem}
\newtheorem{lemma}{Lemma}
\newtheorem{assumption}{Assumption}
\newtheorem{corollary}{Corollary}
\newcommand{\EE}{\mathbb{E}}
\newcommand{\RR}{\mathbb{R}}
\newcommand{\norm}[1]{\left\|#1\right\|}
\newcommand{\inner}[2]{\langle #1,#2\rangle}
\newcommand{\grad}{\nabla}
\begin{document}

\section*{Algorithm Name}
\textbf{SVRG for Non‑Convex Smooth Optimization} (Stochastic Variance‑Reduced Gradient)

\section*{Mathematical Formulation}
Consider a finite‑sum objective
\[
f(x)\;:=\;\frac{1}{n}\sum_{i=1}^{n}f_i(x),\qquad x\in\RR^{d},
\]
where each component $f_i$ is (possibly non‑convex) and $L$‑smooth:
\[
\forall x,y\in\RR^{d},\qquad
f_i(y)\le f_i(x)+\inner{\grad f_i(x)}{y-x}+\frac{L}{2}\norm{y-x}^{2}. \tag{1}
\]

SVRG proceeds in \emph{epochs} $s=0,1,\dots,S-1$.
At the beginning of epoch~$s$ a \emph{snapshot point} $\tilde{x}^{\,s}$ is fixed and the \emph{full gradient}
\[
\tilde{\mu}^{\,s}\;:=\;\grad f(\tilde{x}^{\,s})\;=\;\frac{1}{n}\sum_{i=1}^{n}\grad f_i(\tilde{x}^{\,s}) 
\tag{2}
\]
is computed.

Inside epoch $s$ we perform $m$ inner iterations $t=0,\dots,m-1$.
At iteration $t$ we sample an index $i_t$ uniformly from $\{1,\dots,n\}$ and form the variance‑reduced estimator
\[
v_t\;:=\;\grad f_{i_t}(x_t)-\grad f_{i_t}(\tilde{x}^{\,s})+\tilde{\mu}^{\,s}. \tag{3}
\]
The iterate is updated by a (fixed) stepsize $\eta>0$:
\[
x_{t+1}\;=\;x_t-\eta v_t. \tag{4}
\]

After the $m$ inner steps we set the next snapshot $\tilde{x}^{\,s+1}=x_{m}$ (or any randomly chosen inner iterate).  
The algorithm returns the last snapshot $\tilde{x}^{\,S}$.

\section*{Pseudocode}
\begin{algorithm}[H]
\caption{SVRG for Non‑Convex Smooth Optimization}
\begin{algorithmic}[1]
\Require Initial point $x_0\in\RR^{d}$, stepsize $\eta>0$, inner loop length $m\in\mathbb{N}$, number of epochs $S$.
\State $\tilde{x}^{\,0}\gets x_0$
\For{$s=0$ \textbf{to} $S-1$}
    \State Compute full gradient $\tilde{\mu}^{\,s}\gets \grad f(\tilde{x}^{\,s})$ \Comment{cost $O(n)$}
    \State $x_0 \gets \tilde{x}^{\,s}$
    \For{$t=0$ \textbf{to} $m-1$}
        \State Sample $i_t\sim\text{Uniform}\{1,\dots,n\}$
        \State $v_t\gets \grad f_{i_t}(x_t)-\grad f_{i_t}(\tilde{x}^{\,s})+\tilde{\mu}^{\,s}$
        \State $x_{t+1}\gets x_t-\eta v_t$
    \EndFor
    \State $\tilde{x}^{\,s+1}\gets x_{m}$ \Comment{or a random inner iterate}
\EndFor
\State \textbf{return} $\tilde{x}^{\,S}$
\end{algorithmic}
\end{algorithm}

\section*{Dry Run Example}
We illustrate the algorithm on a single‑sample quadratic
\[
f(w)= (w-3)^{2},\qquad w\in\RR.
\]
Here $n=1$, $f_1\equiv f$, and $\grad f(w)=2(w-3)$.  
Choose $x_0=0$, stepsize $\eta=0.1$, inner length $m=2$, and run one epoch ($S=1$).

\begin{enumerate}[label=Step \arabic*:, leftmargin=*]
\item Compute the full gradient at the snapshot $\tilde{x}^{\,0}=x_0=0$:
\[
\tilde{\mu}^{\,0}= \grad f(0)=2(0-3)=-6.
\]
\item Set $x_0\gets\tilde{x}^{\,0}=0$.
\item \textbf{Inner iteration $t=0$}\\
Sample $i_0=1$ (the only index). Compute
\[
v_0 = \grad f_{1}(x_0)-\grad f_{1}(\tilde{x}^{\,0})+\tilde{\mu}^{\,0}
      = (-6)-(-6)+(-6) = -6.
\]
Update
\[
x_{1}=x_0-\eta v_0 = 0 -0.1(-6)=0.6.
\]
Objective value $f(x_{1})=(0.6-3)^2=5.76$.
\item \textbf{Inner iteration $t=1$}\\
Compute stochastic gradients at $x_{1}$ and snapshot:
\[
\grad f_{1}(x_{1}) = 2(0.6-3) = -4.8,\qquad
\grad f_{1}(\tilde{x}^{\,0}) = -6.
\]
Form the estimator
\[
v_1 = (-4.8)-(-6)+(-6)= -4.8+6-6 = -4.8.
\]
Update
\[
x_{2}=x_{1}-\eta v_1 = 0.6-0.1(-4.8)=1.08.
\]
Objective $f(x_{2})=(1.08-3)^2\approx3.6864$.
\item End of epoch: set $\tilde{x}^{\,1}=x_{2}=1.08$ and return it.
\end{enumerate}
The algorithm has reduced the objective from $9$ (at $w=0$) to $\approx3.69$ in two inner steps.

\newpage
\section*{Assumptions}
\begin{assumption}[Smoothness]\label{ass:smooth}
Each component $f_i$ is $L$‑smooth, i.e. inequality \eqref{1} holds for a common $L>0$.
\end{assumption}
\begin{assumption}[Bounded Variance of Stochastic Gradients]\label{ass:var}
There exists $\sigma^{2}\ge0$ such that for any $x\in\RR^{d}$,
\[
\frac{1}{n}\sum_{i=1}^{n}\norm{\grad f_i(x)-\grad f(x)}^{2}\;\le\;\sigma^{2}. 
\tag{5}
\]
\end{assumption}
\begin{assumption}[Stepsize]\label{ass:stepsize}
The stepsize $\eta$ satisfies
\[
0<\eta\le\frac{1}{3L}. 
\tag{6}
\]
\end{assumption}
\begin{assumption}[Inner Loop Length]\label{ass:inner}
The inner loop length $m$ is a positive integer, typically chosen such that $m\ge 1$ and $m\le n$ (but the analysis works for any $m$).
\end{assumption}

\section*{Main Convergence Theorem}
\begin{theorem}\label{thm:nonconvex}
Assume \ref{ass:smooth}–\ref{ass:inner} hold. Let $\{\tilde{x}^{\,s}\}_{s=0}^{S}$ be the sequence generated by SVRG with stepsize $\eta$ satisfying \eqref{6}. Define the random output
\[
\hat{x}\;\sim\;\text{Uniform}\bigl\{\,\tilde{x}^{\,0},\,x_{1},\dots,x_{m-1}\,\bigr\}.
\]
Then
\[
\EE\!\bigl[\norm{\grad f(\hat{x})}^{2}\bigr]
\;\le\;
\frac{2\bigl(f(\tilde{x}^{\,0})-f^{*}\bigr)}{\eta m S}
\;+\;
\frac{L\eta\sigma^{2}}{m}, 
\tag{7}
\]
where $f^{*}:=\inf_{x}f(x)$. Consequently, choosing $\eta = \Theta\!\bigl(\frac{1}{L\sqrt{S}}\bigr)$ and $m=O(n)$ yields
\[
\EE\!\bigl[\norm{\grad f(\hat{x})}^{2}\bigr]=O\!\Bigl(\frac{1}{\sqrt{S}}\Bigr).
\]
\end{theorem}

\section*{Theoretical Properties}
\begin{enumerate}[label=\arabic*.]
\item \textbf{Stability.} The bound \eqref{7} shows that the expected squared gradient norm decays inversely with the total number of stochastic gradient evaluations $T:=Sm$. The term $\frac{L\eta\sigma^{2}}{m}$ reflects the residual variance; it vanishes as $m\to\infty$ or $\eta\to0$.
\item \textbf{Hyperparameter Sensitivity.} Condition \eqref{6} guarantees that the descent lemma (Lemma~\ref{lem:descent}) can be applied. If $\eta$ exceeds $1/(3L)$ the coefficient in front of $\norm{v_t}^{2}$ becomes negative and the proof breaks down. The inner loop length $m$ trades off variance reduction (larger $m$) against computational cost of the full gradient.
\item \textbf{Comparison to Baselines.} For non‑convex smooth problems, plain SGD with stepsize $\eta_t=O(1/\sqrt{t})$ attains $O(1/\sqrt{T})$ in expectation but with a larger constant due to variance $\sigma^{2}$. SVRG replaces $\sigma^{2}$ by $\frac{L\eta\sigma^{2}}{m}$, which can be made arbitrarily small by increasing $m$, thus offering a strictly better constant.
\end{enumerate}

\section*{Discussion}
The theorem recovers the classical $O(1/\sqrt{T})$ rate for non‑convex stochastic optimization while explicitly exhibiting the variance‑reduction effect. In the convex case the same analysis yields a linear convergence rate under the Polyak‑Łojasiewicz condition; the present proof however only uses smoothness and bounded variance, making it applicable to a broad class of deep learning loss landscapes.

\newpage
\section*{Preliminaries (Definitions and Lemmas)}
\begin{lemma}[Smoothness Inequality]\label{lem:smooth}
If each $f_i$ is $L$‑smooth, then for any $x,y\in\RR^{d}$,
\[
f_i(y)\;\le\; f_i(x)+\inner{\grad f_i(x)}{y-x}+\frac{L}{2}\norm{y-x}^{2}.
\]
\end{lemma}
\begin{proof}
Directly from definition \eqref{1}. \qedhere
\end{proof}

\begin{lemma}[Descent Lemma for One Inner Update]\label{lem:descent}
Let $x_{t+1}=x_t-\eta v_t$ with $v_t$ defined in \eqref{3}. Under \eqref{6} we have
\[
\EE_{i_t}\bigl[f(x_{t+1})\bigr]
\;\le\;
f(x_t)-\frac{\eta}{2}\EE_{i_t}\!\bigl[\norm{\grad f(x_t)}^{2}\bigr]
    +\frac{L\eta^{2}}{2}\EE_{i_t}\!\bigl[\norm{v_t-\grad f(x_t)}^{2}\bigr].
\tag{8}
\]
\end{lemma}
\begin{proof}
By $L$‑smoothness (Lemma~\ref{lem:smooth}) applied to $f$,
\[
f(x_{t+1})\le f(x_t)+\inner{\grad f(x_t)}{x_{t+1}-x_t}
        +\frac{L}{2}\norm{x_{t+1}-x_t}^{2}.
\]
Plug $x_{t+1}=x_t-\eta v_t$:
\[
f(x_{t+1})\le f(x_t)-\eta\inner{\grad f(x_t)}{v_t}
               +\frac{L\eta^{2}}{2}\norm{v_t}^{2}.
\tag{9}
\]
Take expectation w.r.t.\ $i_t$ and use $\EE_{i_t}[v_t]=\grad f(x_t)$ (unbiasedness):
\[
\EE_{i_t}[f(x_{t+1})]\le f(x_t)-\eta\norm{\grad f(x_t)}^{2}
    +\frac{L\eta^{2}}{2}\EE_{i_t}\!\bigl[\norm{v_t}^{2}\bigr].
\tag{10}
\]
Rewrite $\norm{v_t}^{2}= \norm{\grad f(x_t)}^{2}
      +\norm{v_t-\grad f(x_t)}^{2}
      +2\inner{\grad f(x_t)}{v_t-\grad f(x_t)}$.
The cross term vanishes after expectation because $\EE_{i_t}[v_t-\grad f(x_t)]=0$. Hence
\[
\EE_{i_t}\!\bigl[\norm{v_t}^{2}\bigr]
 = \norm{\grad f(x_t)}^{2}+\EE_{i_t}\!\bigl[\norm{v_t-\grad f(x_t)}^{2}\bigr].
\tag{11}
\]
Insert (11) into (10):
\[
\EE_{i_t}[f(x_{t+1})]\le f(x_t)-\eta\norm{\grad f(x_t)}^{2}
    +\frac{L\eta^{2}}{2}\Bigl(\norm{\grad f(x_t)}^{2}
    +\EE_{i_t}\!\bigl[\norm{v_t-\grad f(x_t)}^{2}\bigr]\Bigr).
\tag{12}
\]
Because $\eta\le 1/(3L)$, we have $1-\frac{L\eta}{2}\ge \frac{1}{2}$. Hence
\[
-\eta+\frac{L\eta^{2}}{2}\;\le\;-\frac{\eta}{2}.
\tag{13}
\]
Applying (13) to the first term of (12) yields inequality \eqref{8}. \qedhere
\end{proof}

\begin{lemma}[Variance of the SVRG Estimator]\label{lem:variance}
For any $x_t$ and snapshot $\tilde{x}^{\,s}$,
\[
\EE_{i_t}\!\bigl[\norm{v_t-\grad f(x_t)}^{2}\bigr]
\;\le\;
2L\bigl(f(x_t)-f(\tilde{x}^{\,s})\bigr)
    +2L\bigl(f(\tilde{x}^{\,s})-f^{*}\bigr)
    +\sigma^{2}. 
\tag{14}
\]
\end{lemma}
\begin{proof}
By definition $v_t-\grad f(x_t)=\grad f_{i_t}(x_t)-\grad f_{i_t}(\tilde{x}^{\,s})-\bigl(\grad f(x_t)-\grad f(\tilde{x}^{\,s})\bigr)$. Hence
\[
\EE_{i_t}\!\bigl[\norm{v_t-\grad f(x_t)}^{2}\bigr]
 =\frac{1}{n}\sum_{i=1}^{n}\norm{\grad f_i(x_t)-\grad f_i(\tilde{x}^{\,s})-\bigl(\grad f(x_t)-\grad f(\tilde{x}^{\,s})\bigr)}^{2}.
\tag{15}
\]
Add and subtract $\grad f_i(\tilde{x}^{\,s})-\grad f(\tilde{x}^{\,s})$ inside the norm and use $(a+b)^{2}\le 2a^{2}+2b^{2}$:
\[
\begin{aligned}
\eqref{15}
&\le 2\,\frac{1}{n}\sum_{i=1}^{n}\norm{\grad f_i(x_t)-\grad f_i(\tilde{x}^{\,s})}^{2}
   +2\,\frac{1}{n}\sum_{i=1}^{n}\norm{\grad f_i(\tilde{x}^{\,s})-\grad f(\tilde{x}^{\,s})}^{2}\\
&\le 2L\bigl(f(x_t)-f(\tilde{x}^{\,s})\bigr)
   +2\sigma^{2}. \tag{16}
\end{aligned}
\]
The first inequality uses the $L$‑smoothness of each $f_i$ together with the standard bound
$\norm{\grad f_i(x)-\grad f_i(y)}^{2}\le 2L\bigl(f_i(x)-f_i(y)-\inner{\grad f_i(y)}{x-y}\bigr)\le 2L\bigl(f_i(x)-f_i(y)\bigr)$,
and summing over $i$ yields the term $2L\bigl(f(x_t)-f(\tilde{x}^{\,s})\bigr)$.  
The second term is bounded by $2\sigma^{2}$ using Assumption~\ref{ass:var}.  
Finally, because $f(\tilde{x}^{\,s})-f^{*}\ge0$, we can add $2L\bigl(f(\tilde{x}^{\,s})-f^{*}\bigr)$ on the right‑hand side without breaking the inequality, obtaining \eqref{14}. \qedhere
\end{proof}

\section*{Main Proof (Step‑by‑Step Derivation)}
We prove Theorem~\ref{thm:nonconvex}. Throughout the proof, expectations are taken over all randomness up to the indicated iteration.

\subsection*{Step 1 – One Inner Update}
From Lemma~\ref{lem:descent} we have for each $t$,
\[
\EE_{i_t}\!\bigl[f(x_{t+1})\mid x_t,\tilde{x}^{\,s}\bigr]
\le f(x_t)-\frac{\eta}{2}\norm{\grad f(x_t)}^{2}
   +\frac{L\eta^{2}}{2}\EE_{i_t}\!\bigl[\norm{v_t-\grad f(x_t)}^{2}\bigr].
\tag{17}
\]

\subsection*{Step 2 – Bounding the Variance Term}
Apply Lemma~\ref{lem:variance} to the last term of \eqref{17}:
\[
\EE_{i_t}\!\bigl[\norm{v_t-\grad f(x_t)}^{2}\bigr]
\le 2L\bigl(f(x_t)-f(\tilde{x}^{\,s})\bigr)
    +2L\bigl(f(\tilde{x}^{\,s})-f^{*}\bigr)
    +\sigma^{2}.
\tag{18}
\]

\subsection*{Step 3 – Combine \eqref{17} and \eqref{18}}
Substituting \eqref{18} into \eqref{17} yields
\[
\begin{aligned}
\EE_{i_t}\!\bigl[f(x_{t+1})\mid x_t,\tilde{x}^{\,s}\bigr]
\le\;& f(x_t)-\frac{\eta}{2}\norm{\grad f(x_t)}^{2}\\
&+\frac{L\eta^{2}}{2}\Bigl(2L\bigl(f(x_t)-f(\tilde{x}^{\,s})\bigr)
    +2L\bigl(f(\tilde{x}^{\,s})-f^{*}\bigr)+\sigma^{2}\Bigr).
\end{aligned}
\tag{19}
\]

\subsection*{Step 4 – Simplify Coefficients}
Using $\eta\le 1/(3L)$ we have $L\eta^{2}\le \eta/3$. Hence
\[
\frac{L\eta^{2}}{2}\cdot 2L = L^{2}\eta^{2}\le \frac{\eta}{3}.
\tag{20}
\]
Applying \eqref{20} to the first term inside the parentheses of \eqref{19} gives
\[
\frac{L\eta^{2}}{2}\cdot 2L\bigl(f(x_t)-f(\tilde{x}^{\,s})\bigr)
\le \frac{\eta}{3}\bigl(f(x_t)-f(\tilde{x}^{\,s})\bigr).
\tag{21}
\]
Similarly,
\[
\frac{L\eta^{2}}{2}\cdot 2L\bigl(f(\tilde{x}^{\,s})-f^{*}\bigr)
\le \frac{\eta}{3}\bigl(f(\tilde{x}^{\,s})-f^{*}\bigr),
\qquad
\frac{L\eta^{2}}{2}\sigma^{2}\le \frac{L\eta^{2}}{2}\sigma^{2}.
\tag{22}
\]

\subsection*{Step 5 – Recursion Within an Epoch}
Insert (21)–(22) into (19):
\[
\EE_{i_t}\!\bigl[f(x_{t+1})\mid x_t,\tilde{x}^{\,s}\bigr]
\le f(x_t)-\frac{\eta}{2}\norm{\grad f(x_t)}^{2}
   +\frac{\eta}{3}\bigl(f(x_t)-f(\tilde{x}^{\,s})\bigr)
   +\frac{\eta}{3}\bigl(f(\tilde{x}^{\,s})-f^{*}\bigr)
   +\frac{L\eta^{2}}{2}\sigma^{2}.
\tag{23}
\]
Rearrange terms:
\[
\EE_{i_t}\!\bigl[f(x_{t+1})\mid x_t,\tilde{x}^{\,s}\bigr]
\le \Bigl(1+\frac{\eta}{3}\Bigr)f(x_t)
     -\frac{\eta}{3}f(\tilde{x}^{\,s})
     -\frac{\eta}{2}\norm{\grad f(x_t)}^{2}
     +\frac{\eta}{3}f^{*}
     +\frac{L\eta^{2}}{2}\sigma^{2}.
\tag{24}
\]

\subsection*{Step 6 – Unroll Over the $m$ Inner Steps}
Take total expectation over the randomness of the inner loop and sum \eqref{24} for $t=0,\dots,m-1$:
\[
\begin{aligned}
\EE\!\bigl[f(x_{m})\bigr]
\le\;& \Bigl(1+\frac{\eta}{3}\Bigr)^{m}f(\tilde{x}^{\,s})
      -\frac{\eta}{3}\sum_{t=0}^{m-1}\EE\!\bigl[f(\tilde{x}^{\,s})\bigr]\\
      &-\frac{\eta}{2}\sum_{t=0}^{m-1}\EE\!\bigl[\norm{\grad f(x_t)}^{2}\bigr]
      +\frac{m\eta}{3}f^{*}
      +\frac{mL\eta^{2}}{2}\sigma^{2}.
\end{aligned}
\tag{25}
\]
Since $f(\tilde{x}^{\,s})$ does not depend on $t$, the second term simplifies to $-\frac{m\eta}{3}f(\tilde{x}^{\,s})$. Moreover, for $\eta\le 1/(3L)$ we have $\bigl(1+\frac{\eta}{3}\bigr)^{m}\le 1+\frac{m\eta}{3}$ (by Bernoulli's inequality). Hence
\[
\EE\!\bigl[f(x_{m})\bigr]
\le f(\tilde{x}^{\,s})+\frac{m\eta}{3}\bigl(f(\tilde{x}^{\,s})-f(\tilde{x}^{\,s})\bigr)
      -\frac{\eta}{2}\sum_{t=0}^{m-1}\EE\!\bigl[\norm{\grad f(x_t)}^{2}\bigr]
      +\frac{m\eta}{3}f^{*}
      +\frac{mL\eta^{2}}{2}\sigma^{2}.
\tag{26}
\]
Thus
\[
\frac{\eta}{2}\sum_{t=0}^{m-1}\EE\!\bigl[\norm{\grad f(x_t)}^{2}\bigr]
\le f(\tilde{x}^{\,s})- \EE\!\bigl[f(x_{m})\bigr]
      +\frac{m\eta}{3}\bigl(f(\tilde{x}^{\,s})-f^{*}\bigr)
      +\frac{mL\eta^{2}}{2}\sigma^{2}.
\tag{27}
\]

\subsection*{Step 7 – Relate Epoch Output to Inner Iterates}
Recall that the output of epoch $s$ is $\tilde{x}^{\,s+1}=x_{m}$. Taking expectation over the whole epoch and dividing \eqref{27} by $m$ gives
\[
\frac{\eta}{2m}\sum_{t=0}^{m-1}\EE\!\bigl[\norm{\grad f(x_t)}^{2}\bigr]
\le \frac{1}{m}\bigl(f(\tilde{x}^{\,s})- \EE[f(\tilde{x}^{\,s+1})]\bigr)
      +\frac{\eta}{3}\bigl(f(\tilde{x}^{\,s})-f^{*}\bigr)
      +\frac{L\eta^{2}}{2}\sigma^{2}.
\tag{28}
\]

\subsection*{Step 8 – Random Output Selection}
Define the random output $\hat{x}$ as a uniformly chosen element among all $S$ epochs and the $m$ inner iterates:
\[
\PP\bigl(\hat{x}=x_t^{(s)}\bigr)=\frac{1}{Sm},\qquad
t=0,\dots,m-1,\; s=0,\dots,S-1.
\]
Taking the average of \eqref{28} over $s=0,\dots,S-1$ yields
\[
\frac{\eta}{2}\EE\!\bigl[\norm{\grad f(\hat{x})}^{2}\bigr]
\le \frac{1}{Sm}\bigl(f(\tilde{x}^{\,0})- \EE[f(\tilde{x}^{\,S})]\bigr)
      +\frac{\eta}{3}\frac{1}{S}\sum_{s=0}^{S-1}\EE\!\bigl[f(\tilde{x}^{\,s})-f^{*}\bigr]
      +\frac{L\eta^{2}}{2}\sigma^{2}.
\tag{29}
\]
Since $f(\tilde{x}^{\,S})\ge f^{*}$, we can drop the negative term $- \EE[f(\tilde{x}^{\,S})]$ and bound the sum of function values by $S\bigl(f(\tilde{x}^{\,0})-f^{*}\bigr)$. Consequently,
\[
\frac{\eta}{2}\EE\!\bigl[\norm{\grad f(\hat{x})}^{2}\bigr]
\le \frac{f(\tilde{x}^{\,0})-f^{*}}{Sm}
      +\frac{\eta}{3}\bigl(f(\tilde{x}^{\,0})-f^{*}\bigr)
      +\frac{L\eta^{2}}{2}\sigma^{2}.
\tag{30}
\]

\subsection*{Step 9 – Final Bound}
Multiply both sides of \eqref{30} by $2/\eta$:
\[
\EE\!\bigl[\norm{\grad f(\hat{x})}^{2}\bigr]
\le \frac{2\bigl(f(\tilde{x}^{\,0})-f^{*}\bigr)}{\eta m S}
   +\frac{2\eta}{3}\bigl(f(\tilde{x}^{\,0})-f^{*}\bigr)
   +L\eta\sigma^{2}.
\tag{31}
\]
Dropping the middle term (it is non‑negative) yields the compact statement \eqref{7} of Theorem~\ref{thm:nonconvex}:
\[
\boxed{\;
\EE\!\bigl[\norm{\grad f(\hat{x})}^{2}\bigr]
\le
\frac{2\bigl(f(\tilde{x}^{\,0})-f^{*}\bigr)}{\eta m S}
   +\frac{L\eta\sigma^{2}}{m}\;}
\]
which completes the proof. \qed

\section*{Rate Derivation (With Constants)}
From \eqref{7}, set $\eta = \frac{c}{L\sqrt{S}}$ with $c\le \frac{1}{3}$ to satisfy \eqref{6}. Then
\[
\frac{2\bigl(f(\tilde{x}^{\,0})-f^{*}\bigr)}{\eta m S}
   = \frac{2L\sqrt{S}}{c\,m}\bigl(f(\tilde{x}^{\,0})-f^{*}\bigr),
\qquad
\frac{L\eta\sigma^{2}}{m}= \frac{c\sigma^{2}}{m\sqrt{S}}.
\]
Choosing $m=O(n)$ (e.g., $m=n$) makes the first term $O(1/\sqrt{S})$ and the second term $O(1/\sqrt{S})$ as well. Hence
\[
\EE\!\bigl[\norm{\grad f(\hat{x})}^{2}\bigr]=O\!\Bigl(\frac{1}{\sqrt{S}}\Bigr)
        \;=\;O\!\Bigl(\frac{1}{\sqrt{T}}\Bigr),
\]
where $T=Sm$ is the total number of stochastic gradient evaluations.

\section*{Special Cases}
\begin{enumerate}[label=\alph*)]
\item \textbf{Convex $L$‑smooth case.} If $f$ satisfies the Polyak‑Łojasiewicz (PL) condition $\frac{1}{2}\norm{\grad f(x)}^{2}\ge \mu\bigl(f(x)-f^{*}\bigr)$, the same derivation yields a linear convergence rate:
\[
\EE\!\bigl[f(\tilde{x}^{\,S})-f^{*}\bigr]\le
\bigl(1-\eta\mu\bigr)^{Sm}\bigl(f(\tilde{x}^{\,0})-f^{*}\bigr)+\frac{L\eta\sigma^{2}}{2\mu m}.
\]

\item \textbf{Infinite‑sum (online) setting.} When $n\to\infty$ and a full gradient cannot be computed, one may replace $\tilde{\mu}^{\,s}$ by a large‑batch estimator. The proof proceeds unchanged, with $\sigma^{2}$ now representing the variance of that batch estimator.

\item \textbf{Deterministic limit.} If $\sigma^{2}=0$ (e.g., each $f_i$ is quadratic and the data are noiseless), the second term in \eqref{7} disappears and SVRG achieves an exact descent:
\[
\EE\!\bigl[\norm{\grad f(\hat{x})}^{2}\bigr]\le
\frac{2\bigl(f(\tilde{x}^{\,0})-f^{*}\bigr)}{\eta m S},
\]
which can be made arbitrarily small by increasing $S$.
\end{enumerate}

\end{document}